% answers_short.tex —— 1998 年·上海卷·理工农医类 · 简短答案速查
% 由 tools/make_short_answers.py 自 answers.tex 的 \gksolhead/\gkanswer 提取，
% 并与 1998/1998_数学真题_上海理_答案要点.md 逐题比对后生成，勿手改（改源 MD 后重跑）。

\begin{gknotice}
\textbf{使用说明}：本节只给最终答案便于快速核对，完整推导见解析版。答案为编写者独立推导，\textbf{非官方标准答案}。
\end{gknotice}

\gkpart{填空题}
\begin{itemize}[leftmargin=2.6em]
  \item[\textbf{1.}] $2$。
  \item[\textbf{2.}] $5$。
  \item[\textbf{3.}] $4$。
  \item[\textbf{4.}] $(x-2)^3+1$（$x\in\mathbb R$）。
  \item[\textbf{5.}] $\dfrac{2\sqrt2}{3}$。
  \item[\textbf{6.}] $\dfrac{x^2}{3}+\dfrac{y^2}{4}=1$。
  \item[\textbf{7.}] $\dfrac{18}{35}$。
  \item[\textbf{8.}] $4$。
  \item[\textbf{9.}] $4$。
  \item[\textbf{10.}] $2$。
  \item[\textbf{11.}] $\dfrac12$ 或 $\dfrac32$。
\end{itemize}

\gkpart{选择题}
\begin{center}\begin{tabular}{cccccc}
\toprule
题号 & 12 & 13 & 14 & 15 & 16 \\
\midrule
答案 & B & A & C & D & C \\
\bottomrule
\end{tabular}\end{center}

\gkpart{解答题}
\begin{itemize}[leftmargin=2.6em]
  \item[\textbf{17.}] $\dfrac{7+24\sqrt3}{50}$。
  \item[\textbf{18.}] $z=3-i$；$z'=\sqrt2-2\sqrt2 i$，$z'+\sqrt2 i=\sqrt2-\sqrt2 i$；按卷内辐角区间 $[0,2\pi)$，$\theta=\dfrac{7\pi}{4}$（若取 $(-\pi,\pi]$ 则写作 $-\dfrac\pi4$）。
  \item[\textbf{19.}] （1）二面角为 $\arctan(2\sqrt3)$；（2）异面直线 $OB_1$ 与 $D_1E$ 所成角为 $\arccos\dfrac{7\sqrt{10}}{40}$。
  \item[\textbf{20.}] 轨迹 $C:\ y^2=4(x-1)$；倾斜角 $\alpha=\dfrac\pi6$ 或 $\dfrac{5\pi}6$。
  \item[\textbf{21.}] $T(t)=t^3-3t+60$；11:00 与 14:00 温度最高，均为 $62^\circ\mathrm C$；8:00 至 16:00 平均温度 $60^\circ\mathrm C$。
  \item[\textbf{22.}] （1）$b_n=-\dfrac{12n+5}{4}$；（2）$\dfrac13$；（3）$c_n=7-24n$（$n\in\mathbb N$）。
\end{itemize}
